\documentclass[11pt]{article}

\usepackage{kotex}
\usepackage{amsmath, amssymb, amsthm}
\usepackage{hyperref}
\usepackage{geometry}
\geometry{margin=1in}

\theoremstyle{definition}
\newtheorem{definition}{Definition}[section]
\newtheorem{remark}[definition]{Remark}
\theoremstyle{plain}
\newtheorem{theorem}[definition]{Theorem}
\newtheorem{lemma}[definition]{Lemma}
\newtheorem{corollary}[definition]{Corollary}

\newcommand{\MSO}{\mathrm{MSO}}
\newcommand{\MSOtwo}{\mathrm{MSO}_2}
\newcommand{\MSOone}{\mathrm{MSO}_1}
\newcommand{\lean}[1]{\texttt{#1}}
\newcommand{\leanfile}[1]{\texttt{#1}}
\newcommand{\BAGS}{\mathrm{BAGS}}
\newcommand{\adh}{\mathrm{adh}}

\title{GraphMSO Lean 형식화 해설 III: \\ MSO Formula over \texorpdfstring{$\tau_P$}{tau P} Graph}
\author{}
\date{}

\begin{document}
\maketitle

\tableofcontents

\section{개요: 세 번의 translation}
\label{mso:sec:overview}

이 문서는 graph 위의 $\MSOtwo$ formula를 최종적으로 tree 위의 tree language
formula로 옮기는 과정을 다룬다. 참조 문서는
\leanfile{Courcelle/lecture\_note\_expanded.tex}이며, 그 Theorem
\lean{thm:courcelle}가 목표다.

전체 경로는 세 단계다.

\begin{center}
\begin{tabular}{clcl}
  (1) & $\MSOtwo$ over $G$ & $\longrightarrow$ & $\MSOone$ over $\tau_I$-incidence structure \\
      & \multicolumn{3}{l}{\quad \S\ref{mso:sec:incidence}, \lean{Formula.toIncidence}} \\[4pt]
  (2) & $\tau_P$-graph $+$ colored decomposition & $\longrightarrow$ & legal $\Sigma$-tree \\
      & \multicolumn{3}{l}{\quad \S\ref{mso:sec:encoding}, \lean{encode}} \\[4pt]
  (3) & $\MSOone$ over $\tau_P$ & $\longrightarrow$ & tree $\MSO$ over $\Sigma_\omega$ \\
      & \multicolumn{3}{l}{\quad \S\ref{mso:sec:translation}, \lean{Formula.translate}} \\
\end{tabular}
\end{center}

(1)은 두 sort(vertex와 edge)를 하나로 합치는 단계이고, (3)이 진짜 Courcelle
translation이다. (2)는 (3)이 말할 대상, 즉 tree model을 만들어 주는 단계다.

이 문서는 pipeline이 실행되는 순서대로 서술한다: 먼저 \S\ref{mso:sec:incidence}
에서 (1)을 다루고, \S\ref{mso:sec:rooted}--\S\ref{mso:sec:decoding}에서 (2)를,
\S\ref{mso:sec:definingpairs}--\S\ref{mso:sec:translation}에서 (3)을 다룬다.
(Lean 파일의 \lean{import} 순서로는 (2)(3)이 먼저이고 (1)이 독립 가지이지만,
이해 순서로는 이쪽이 낫다.)

기초 개념 --- $\tau_P$-graph, tree decomposition, $\BAGS(v)$, bag coloring,
tree language의 syntax와 semantics --- 는 모두
\leanfile{Lean\_Basics.tex}에 있다.

\section{(1) \texorpdfstring{$\MSOtwo$}{MSO2}에서 incidence structure 위의 \texorpdfstring{$\MSOone$}{MSO1}로}
\label{mso:sec:incidence}

\emph{파일: \leanfile{GraphMSO/Decomp/incidenceDecomp.lean},
\leanfile{GraphMSO/incidenceTranslation.lean}}

\emph{참조: \leanfile{lecture\_note\_expanded.tex} Corollary
($\MSOtwo$ properties), 176--230행 부근}

\subsection{Incidence graph의 tree decomposition}

\begin{definition}[\lean{incidenceDecomposition}]
\label{mso:def:incidencedecomp}
$G$의 tree decomposition $D$로부터
$\lean{IncidenceGraph}\ G$ (\leanfile{Lean\_Basics.tex}의
\lean{basics:def:incidencegraph}) 의 tree decomposition을 만든다. 새 tree는
$D$의 tree에 \emph{edge마다 leaf 하나}를 붙인 pendant extension
(\leanfile{Lean\_Basics.tex}의 \lean{basics:def:pendant})이다.
\begin{itemize}
  \item 기존 node $t$의 bag: $\{\lean{fromV}\ v : v \in B_t\}$;
  \item edge $e$의 새 leaf bag: $\{\lean{fromEdge}\ e\} \cup \{\lean{fromV}\ v : v \in e\}$,
    즉 edge object와 그 양 endpoint.
\end{itemize}
새 leaf는 \lean{attachNode e}, 즉 $e$의 두 endpoint를 함께 담는 어떤 node
아래에 붙인다 (\lean{exists\_forall\_mem\_bag}로 존재를 보장하고 선택한다).
\end{definition}

\begin{theorem}[\lean{incidenceDecomposition\_hasWidth}]
\label{mso:thm:incidencewidth}
$D$의 width가 $\omega$ 이하이면 \lean{incidenceDecomposition}의 width는
$\max(\omega, 2)$ 이하다.
\end{theorem}

\begin{remark}[왜 $\max(\omega,2)$인가]
기존 bag은 크기가 그대로이므로 $\omega+1$ 이하다. 새 leaf bag은 원소가
정확히 세 개 (edge object 하나 $+$ endpoint 둘)이므로 $3 = 2+1$ 이하다.
따라서 $\max(\omega,2)$. Edge object의 개수가 아무리 많아도 width는 늘지
않는다는 것이 이 구성의 핵심이다.

$\max(\omega,2)$의 $2$가 나중에 alphabet 크기에 영향을 주므로, $\omega \ge 2$인
경우 --- 즉 실질적으로 관심 있는 모든 경우 --- 에는 width가 전혀 늘지 않는다.
\end{remark}

Running intersection 공리 (\lean{incidence\_connectivity})가 이 구성에서
가장 손이 많이 가는 부분이다. \lean{fromV v}의 occurrence set은 원래
$\BAGS(v)$에 ``$v$를 endpoint로 갖는 edge들의 leaf''를 붙인 것인데, 그 leaf들이
모두 $\BAGS(v)$에 속한 node에 붙어 있으므로 연결성이 유지된다.
\lean{fromEdge e}의 occurrence set은 leaf 하나뿐이라 자명하다.

또 \lean{card\_incidenceVertex\_le\_of\_hasWidth}: incidence structure의 크기가
$(\omega+1)|V|$ 이하다. \leanfile{Lean\_Basics.tex}의 edge bound
\lean{basics:thm:edgebound}에서 나온다.

\subsection{Formula의 translation}

\begin{definition}[\lean{Formula.toIncidence}]
\label{mso:def:toincidence}
$\MSOtwo$ formula를 $\tau_I = \{\mathrm{adj}, \mathrm{Vert}, \mathrm{EdgeObj}\}$
위의 $\MSOone$ formula로 옮긴다. 변수는 \emph{parity}로 합친다:
\[
  \text{vertex 변수 } x \mapsto 2x, \qquad
  \text{edge 변수 } e \mapsto 2e+1,
\]
set 변수도 마찬가지. Atom은:
\[
\begin{array}{lcl}
  \lean{equal}\ x\ y &\mapsto& \lean{equal}\ (2x)\ (2y) \\
  \lean{inSet}\ x\ X &\mapsto& \lean{inSet}\ (2x)\ (2X) \\
  \lean{equalEdge}\ e\ f &\mapsto& \lean{equal}\ (2e{+}1)\ (2f{+}1) \\
  \lean{inc}\ x\ e &\mapsto& \lean{adj}\ (2x)\ (2e{+}1) \\
  \lean{inEdgeSet}\ e\ F &\mapsto& \lean{inSet}\ (2e{+}1)\ (2F{+}1) \\
  \lean{edge}\ x\ y &\mapsto& \exists z.\ \mathrm{EdgeObj}(z) \land \mathrm{adj}(2x, z) \land \mathrm{adj}(2y, z) \land 2x \ne 2y
\end{array}
\]
Quantifier는 sort predicate로 guard한다: $\exists x\,\varphi \mapsto \exists (2x).\,\mathrm{Vert}(2x) \land \varphi'$,
$\forall x\,\varphi \mapsto \forall (2x).\,\mathrm{Vert}(2x) \to \varphi'$, 등.
Set quantifier는 pointwise guard \lean{vertSetGuard} / \lean{edgeSetGuard}
(``$X$의 모든 원소가 \lean{Vert}이다'')를 쓴다.
\end{definition}

\begin{remark}[\lean{edge x y} 번역의 세부사항 두 가지]
\label{mso:rem:edgeatom}
$\tau_I$에는 원래 graph의 adjacency가 없으므로 (incidence graph에서 원래
vertex끼리는 절대 인접하지 않는다) ``$x$와 $y$를 잇는 edge object가 존재한다''로
표현해야 한다. 여기에 두 가지 기술적 요점이 있다.

\emph{(a) $x \ne y$ 조건이 필수다.} 이것이 없으면 $x = y$이고 $x$에 어떤 edge가
붙어 있을 때 $\lean{edge}\ x\ x$가 참이 되어 버린다. 원래 graph는 simple
graph이므로 loop가 없어야 하는데, incidence 조건 ``$x$가 $z$에 incident''를 두 번
쓰는 것만으로는 두 endpoint가 다르다는 것을 강제하지 못한다.

\emph{(b) bound 변수를 $2x + 2y + 1$로 잡았다.} 이 값은 \emph{홀수}이므로
$2x$, $2y$와 절대 겹치지 않는다. 즉 parity 규약이 자동으로 capture를 막아
준다. \leanfile{Lean\_Basics.tex}의 \lean{basics:rem:subsetcapture}에서 말한
``bound 변수를 구체적으로 고르고 그 선택이 옳음을 확인한다'' 방식의 전형적인
예다. (한편 \lean{vertSetGuard}는 bound 변수로 $0$을 쓰는데, $0 = 2 \cdot 0$은
짝수라서 vertex 변수 $0$과 충돌할 수 있다. 실제로는 guard가 quantifier 바로
안쪽에서만 쓰이고 그 자리에서 자유 변수 $0$이 살아 있지 않도록 정리 진술이
관리한다.)
\end{remark}

\begin{theorem}[\lean{satisfiesAt\_toIncidence\_iff}]
\label{mso:thm:toincidence}
자유 변수 위에서 대응하는 assignment 쌍에 대해, translation의 satisfaction은
원래 $\MSOtwo$ satisfaction과 같다.
\end{theorem}

\begin{remark}[대응 조건이 정리 진술의 대부분이다]
이 정리의 hypothesis는 ``$\MSOtwo$ assignment $\rho$와 $\MSOone$ assignment
$\sigma$가 자유 변수에서 대응한다''인데, 그 대응이 네 종류 변수마다 다른
모양이다. 예를 들어 vertex set 변수 $X$에 대해서는
$\sigma(2X) = \{\lean{fromV}\ v : v \in \rho(X)\}$이어야 한다.
증명은 formula에 대한 structural induction이며, set quantifier 경우에
``guard를 만족하는 $\MSOone$ 쪽 집합은 반드시 \lean{fromV}의 상이다''
(\lean{eq\_image\_fromV\_of\_forall\_isVertex}) 를 써서 역방향 witness를
만든다. 이것이 sort predicate가 필요한 이유
(\leanfile{Lean\_Basics.tex}의 \lean{basics:rem:whysorts})의 구체적 발현이다.
\end{remark}

\begin{corollary}[\lean{satisfies\_toIncidence\_iff}, \lean{incidence\_reduction}]
\label{mso:cor:incidencereduction}
Sentence $\varphi$에 대해
$G \models \varphi \iff \lean{incidenceTauGraph}\ G \models \varphi^{\mathrm{inc}}$.
따라서 width $\omega$ decomposition을 가진 $G$ 위의 $\MSOtwo$ model checking은
width $\max(\omega,2)$ decomposition을 가진 incidence structure 위의 $\MSOone$
model checking으로 환원된다.
\end{corollary}

이로써 이후 모든 논의를 $\tau_P$-graph 위의 $\MSOone$으로 좁힐 수 있다.

\section{(2a) Rooted graph와 gluing}
\label{mso:sec:rooted}

\emph{파일: \leanfile{GraphMSO/Decomp/rootedGraph.lean}}

\emph{참조: \leanfile{lecture\_note\_expanded.tex}
\S``Rooted unary-expanded graphs and gluing''}

\begin{definition}[\lean{KRootedGraph}]
\label{mso:def:krootedgraph}
$\tau_P$-graph에 boundary $R \subseteq V$와 부분 injective labeling
$\lean{label} : \lean{labelDom} \to \lean{Fin}\ k$ (정의역이 $R$을 포함)를
추가한 structure. 참조 문서의 ``$k$-rooted $\tau_P$-graph $(H,R,\iota)$''다.
\end{definition}

Gluing은 \emph{부분적} 연산이다: 두 rooted graph는 두 번째의 boundary를
label로 첫 번째 안에서 찾을 수 있을 때만 붙일 수 있다. Lean 형식화는 이것을
부분 함수로 정의하는 대신 세 층으로 나눈다.

\begin{definition}[gluing의 세 층]
\label{mso:def:gluinglayers}
\begin{itemize}
  \item \lean{Gluable A B} --- 붙일 수 있다는 \emph{전제조건}
    (\lean{LabelCompatibleOnRoots}: $B$의 root label이 모두 $A$에 나타나고
    label이 일치).
  \item \lean{GluingData A B C} --- ``$C$가 $A$와 $B$를 붙인 결과다''라는
    명시적 \emph{witness 데이터}: $A \oplus B$에서 $C$로 가는 map들과 그것들이
    만족하는 성질들.
  \item \lean{IsGluing A B C} $:= \lean{Nonempty}\ (\lean{GluingData}\ A\ B\ C)$
    --- 정리 진술에 쓰는 \emph{명제} 층.
  \item \lean{MultiGluingData}, \lean{IsMultiGluing} --- 하나의 $A$에 family
    $B : \iota \to \lean{KRootedGraph}$를 한꺼번에 붙이는 판.
\end{itemize}
\end{definition}

\begin{remark}[왜 세 층으로 나누는가]
\label{mso:rem:whythreelayers}
부분 함수를 Lean에서 직접 정의하려면 ``정의되지 않은 경우''의 값을 정하거나
전제조건을 인자로 받아야 하는데, 후자는 그 전제조건의 증명이 결과값에 딸려
다니게 되어 두 gluing 결과를 비교할 때 증명 무관성(proof irrelevance) 문제가
생긴다.

세 층 구조는 이 문제를 피한다. 정리는 proof-irrelevant한 \lean{IsGluing}으로
쓰고, 증명 중에 구체적인 map이 필요하면 \lean{GluingData}를 꺼내 쓴다.
결과적으로 ``gluing한 것이 유일하다''를 억지로 증명하지 않아도 되고,
동형인 여러 결과를 모두 받아들일 수 있다. \leanfile{CLAUDE.md}의 지침
(불필요한 \lean{Equiv}를 만들지 말 것)과 일치하는 설계다.
\end{remark}

\section{(2b) Node graph와 cone graph}
\label{mso:sec:nodecone}

\emph{파일: \leanfile{GraphMSO/Decomp/nodeConeGraph.lean}}

\emph{참조: \leanfile{lecture\_note\_expanded.tex} Lemma \lean{lem:cone-struct},
Lemma \lean{lem:cone-gluing}}

\begin{definition}[\lean{nodeGraph}, \lean{coneGraph}]
\label{mso:def:nodeconegraph}
Decomposition의 node $t$에서:
\begin{itemize}
  \item \lean{nodeGraph}: 현재 bag $B_t$가 유도하는 graph. Boundary는 adhesion
    $\adh(t)$, labeling은 bag 전체 위의 coloring.
  \item \lean{coneGraph}: cone $\lean{cone}(t)$ (부분트리 전체의 vertex)가
    유도하는 graph. Boundary는 여전히 $\adh(t)$뿐이지만 labeling은 현재 bag
    위에 남겨 둔다.
\end{itemize}
\end{definition}

\begin{remark}[labeling을 bag 전체에 남기는 이유]
Cone graph의 boundary는 $\adh(t)$지만, label의 정의역은 $B_t$ 전체다.
이렇게 해야 ``node graph를 자식 cone들과 붙인 결과''와 ``부모의 cone graph''가
\emph{같은 labeled interface}를 갖게 되어 두 rooted graph를 직접 비교할 수
있다. Boundary와 label domain을 분리해 둔 \lean{KRootedGraph}의 설계
(\S\ref{mso:def:krootedgraph})가 여기서 쓰인다.
\end{remark}

구조 보조정리들이 참조 문서 Lemma \lean{lem:cone-struct}에 대응한다.
\begin{itemize}
  \item \lean{cone\_inter\_bag\_eq\_adhesion\_of\_isChild}:
    자식 $s$에 대해 $\lean{cone}(s) \cap B_t = \adh(s)$.
  \item \lean{cone\_inter\_cone\_subset\_adhesion\_inter\_adhesion\_of\_isChild\_ne}:
    서로 다른 두 자식의 cone은 두 adhesion 안에서만 만난다.
  \item \lean{adj\_in\_bag\_or\_child\_cone\_of\_adj\_of\_mem\_cone}: cone 안의
    edge는 현재 bag 안에 있거나 어떤 자식 cone 안에 있다.
\end{itemize}

\begin{theorem}[\lean{cone\_isMultiGluing}, 참조 문서 Lemma \lean{lem:cone-gluing}]
\label{mso:thm:conegluing}
Node graph를 모든 자식 cone과 붙인 것은 정확히 부모의 cone graph다. 구체적인
\lean{MultiGluingData} witness가 제시된다.
\end{theorem}

\begin{remark}
이 정리가 ``decomposition을 tree 모양으로 조립할 수 있다''는 사실의 형식화이며,
아래 decoding (\S\ref{mso:sec:decoding})이 역방향으로 하는 일의 근거다.
증명은 gluing 관계 \lean{coneGluingRepr}가 surjective임
(\lean{coneGluingRepr\_surjective}) 과 그 동일시 조건이 정확함
(\lean{coneGluingRepr\_eq\_iff}) 을 보이는 것으로 이루어진다.
\end{remark}

\section{(2c) \texorpdfstring{$\Sigma$}{Sigma}-letter와 \texorpdfstring{$\Sigma$}{Sigma}-tree}
\label{mso:sec:sigmatree}

\emph{파일: \leanfile{GraphMSO/Decomp/sigmaTree.lean}}

\emph{참조: \leanfile{lecture\_note\_expanded.tex} \S``The alphabet and legal
trees'', Definition \lean{def:sigma-tree}}

\begin{definition}[\lean{SigmaLetter}]
\label{mso:def:sigmaletter}
고정된 $\omega$에 대해 $\Sigma_\omega$의 letter는 네 field를 갖는 structure다:
\begin{itemize}
  \item \lean{verts} $\subseteq \{0,\dots,\omega\}$ --- 나타나는 \emph{색}들;
  \item \lean{G} --- \lean{verts} 위의 simple graph;
  \item \lean{R} $\subseteq$ \lean{verts} --- boundary;
  \item \lean{tag} $: P \to \lean{verts} \to \lean{Prop}$ --- 각 색에 대해
    어떤 unary predicate가 성립하는지.
\end{itemize}
\end{definition}

\begin{remark}[letter의 vertex는 \emph{색 자체}다]
\label{mso:rem:lettercolors}
이 점을 놓치면 이후 전체가 헷갈린다. \S\ref{mso:def:krootedgraph}의
\lean{KRootedGraph}에서는 vertex가 원래 graph의 vertex이고 label이 따로
붙어 있다. \lean{SigmaLetter}에서는 vertex \emph{가} 색이다. 즉 encoding은
bag의 vertex를 그 색으로 \emph{개명}하는 작업이며, bag coloring이
bag 위에서 injective하기 때문에 이 개명이 bag 안에서는 정보를 잃지 않는다.

이 개명이 결정적인 이유는, 그래야 alphabet이 \emph{원래 graph와 무관한}
유한 집합이 되기 때문이다. \lean{SigmaLetter.instFinite}가 $P$가 유한하면
$\lean{SigmaLetter}\ P\ \omega$도 유한함을 증명한다 ($\{0,\dots,\omega\}$의
부분집합, 그 위의 graph, boundary, tag가 모두 유한한 선택지이므로).
Tree automaton은 유한 alphabet을 요구하므로 이것 없이는
\leanfile{Lean\_Automata.tex}로 넘어갈 수 없다.

\lean{toKRootedGraph}는 letter를 다시 \lean{KRootedGraph}로 보는 함수인데,
labeling으로 ``색을 자기 자신에 보내는'' identity를 쓴다.
\end{remark}

Letter를 읽는 술어들: \lean{HasVertex} $i$ (색 $i$가 나타난다),
\lean{RootContains} $i$ (색 $i$가 boundary에 있다), \lean{RootEmpty},
\lean{AdjOnColors} $i\ j$ (색 $i$와 $j$가 인접), \lean{TagOnColor} $p\ i$
(색 $i$가 술어 $p$를 만족).

\begin{definition}[\lean{Compatible}, legality]
\label{mso:def:compatible}
child letter가 parent letter와 \lean{Compatible}하다는 것은 세 조건이다:
child의 root 색이 모두 parent에 나타나고, 그 root 색들 사이의 adjacency가
양쪽에서 일치하며, 그 root 색들의 unary predicate tag가 양쪽에서 일치한다.
참조 문서의 $R_x \subseteq V(H_y)$ 및 $H_x[R_x] = H_y[R_x]$에 해당한다.

\lean{SigmaTree}는 tree 구조와 root, 그리고 node마다의 letter를 갖는
structure이고, \lean{IsLegal}은 ``root letter의 boundary가 비어 있고, 모든
non-root node의 letter가 부모와 compatible''이다.
\end{definition}

\begin{remark}[unordered tree다]
\label{mso:rem:unordered}
참조 문서 Definition \lean{def:sigma-tree}는 $\Sigma$-tree를 \emph{순서 있는}
이진 트리 (\lean{child₁}, \lean{child₂})로 정의한다. Lean의 \lean{SigmaTree}는
\leanfile{Lean\_Basics.tex}의 tree decomposition과 같은 방식, 즉
\lean{SimpleGraph.IsTree}와 root로 표현된 \emph{순서 없는} tree다.

이 차이는 의도적이며 \leanfile{Lean\_Basics.tex}의
\lean{basics:rem:treesimplify}에서 설명한 것과 같은 이유다: translation 층은
순서를 전혀 쓰지 않는다. 순서는 automaton에 넘길 때 비로소 필요해지고,
그 지점에서 \leanfile{Lean\_Automata.tex}의 \lean{orderedEncode}와
\lean{orderedEncodeIso}가 순서 있는 표현으로 옮긴다. 두 표현을 잇는 것이
\leanfile{Lean\_Basics.tex}의 isomorphism invariance
(\lean{basics:thm:isoinvariance})다.
\end{remark}

\section{(2d) Encoding}
\label{mso:sec:encoding}

\emph{파일: \leanfile{GraphMSO/Decomp/encoding.lean}}

\begin{definition}[\lean{encodeLetter}]
\label{mso:def:encodeletter}
$\tau_P$-graph $G$, bag-injective coloring $c$, rooted decomposition $T$가
주어졌을 때 node $t$의 letter는:
\[
\begin{array}{ll}
  \lean{verts} &= c(B_t) \quad (\text{bag의 색들}) \\
  \lean{R} &= c(\adh(t)) \quad (\text{adhesion의 색들}) \\
  \lean{G}.\lean{Adj}\ i\ j &\iff \exists u,v \in B_t.\ c(u)=i \land c(v)=j \land u \sim_G v \\
  \lean{tag}\ p\ i &\iff \exists u \in B_t.\ c(u)=i \land \lean{vpred}\ p\ u
\end{array}
\]
\lean{encode}는 같은 tree와 root에 이 letter를 붙인 \lean{SigmaTree}다.
\end{definition}

\begin{remark}[존재 한정으로 쓴 이유]
Adjacency와 tag를 ``색 $i$를 갖는 \emph{그} vertex''로 쓰지 않고 존재
한정으로 쓴 점에 주목하라. 색 $i$를 갖는 bag vertex는 bag-injectivity에 의해
유일하지만, 그 유일성을 이용해 함수를 정의하려면 선택을 해야 하고 그러면
정의가 noncomputable해지며 증명도 무거워진다. 존재 한정으로 쓰면 정의는
가볍고, 유일성은 필요한 자리에서 \lean{hcolor}를 적용해 꺼내 쓴다.

실제로 \lean{loopless} 증명이 그 예다: $i \sim i$이면 $c(u)=c(v)=i$인
$u,v \in B_t$가 있고 bag-injectivity로 $u=v$인데, 원래 graph에 loop가 없으므로
모순이다.
\end{remark}

\begin{theorem}[\lean{encode\_isLegal}]
\label{mso:thm:encodeislegal}
모든 encoding은 legal하다. (참조 문서 Lemma \lean{lem:legal-iff-enc}의
정방향.)
\end{theorem}

\begin{remark}[증명 방식]
Root letter의 boundary가 비어 있는 것은 $\adh(\lean{root}) = \emptyset$이므로
바로 나온다. Compatibility (\lean{encodeLetter\_compatible})가 본론인데,
핵심 논증은 다음과 같다: root 색 $i$는 adhesion의 어떤 vertex $u_0$를
지목하고, adhesion은 $B_t$와 $B_{p(t)}$ 양쪽에 들어 있으므로 bag-injectivity에
의해 $u_0$는 \emph{양쪽 bag에서 색 $i$를 갖는 유일한 vertex}다. 따라서 두
letter가 색 $i$에 대해 읽는 데이터는 같은 vertex에서 나온 것이고, 자동으로
일치한다.

이 정리는 실제 사용처에서 훨씬 일반적인 형태로 서술되어 있다:
\lean{encodeLetter\_compatible}는 $t$와 $s$가 부모-자식 관계라고 가정하지 않고
$\adh(t) = B_t \cap B_s$만 가정한다. \leanfile{CLAUDE.md}의 ``정리의 범용성
유지'' 지침에 맞는 형태다.
\end{remark}

\begin{corollary}[\lean{exists\_encode\_isLegal}]
\label{mso:cor:existsencode}
Width $\omega$ 이하의 decomposition만 있으면 bag coloring이 존재하고
(\leanfile{Lean\_Basics.tex}의 \lean{basics:thm:existsbagcoloring}) 그 encoding이
legal하다. Constructor-coded 판
(\lean{InductiveNiceTreeDecomposition.exists\_encode\_isLegal}) 도 함께
제공되며, \leanfile{Lean\_Nice.tex}의 realization 전달
(\lean{nice:thm:transfer})을 두 번 건너 얻는다.
\end{corollary}

\section{(2e) Decoding}
\label{mso:sec:decoding}

\emph{파일: \leanfile{GraphMSO/Decomp/decoding.lean} (1015줄)}

\emph{참조: \leanfile{lecture\_note\_expanded.tex} \S``Decoding a legal tree'',
Lemma \lean{lem:decode-enc}}

이 절은 encoding의 \emph{역방향}이다: legal $\Sigma$-tree가 주어지면 그것이
어떤 $\tau_P$-graph의 encoding임을 보인다.

\begin{definition}[decoded graph의 quotient 구성]
\label{mso:def:decodedvertex}
\begin{itemize}
  \item \lean{Occurrence} $S := \Sigma\, t : S.\lean{Node},\ (S.\lean{letter}\ t).\lean{verts}$
    --- ``node $t$에서의 색 $i$'' 쌍, 즉 참조 문서의 slot $v(t,i)$.
  \item \lean{colorGraph} $i$ --- 색 $i$를 갖는 node들 위의 graph. $x \sim y$는
    ``tree에서 인접하고, 아래쪽 node의 letter가 $i$를 boundary에 갖는다''.
    즉 \emph{$i$가 그 edge를 따라 실제로 공유된다}는 뜻.
  \item \lean{OccRel} $a\ b$ --- 두 slot이 같은 색 $i$를 갖고,
    \lean{colorGraph} $i$에서 서로 도달 가능하다.
  \item \lean{DecodedVertex} $:= \lean{Quotient}$ of \lean{OccRel}.
\end{itemize}
그 위에 \lean{decodedAdj}, \lean{decodedPred}로 graph와 술어를 정의해
\lean{decodeTauGraph}를 얻고, \lean{decodedBag}, \lean{decodeDecomposition}로
decomposition을 복원한다.
\end{definition}

\begin{remark}[quotient가 gluing의 형식화다]
\label{mso:rem:quotientgluing}
\S\ref{mso:thm:conegluing}의 gluing을 실제로 수행하는 방법이 이 quotient다.
각 node가 자기 letter만큼의 slot을 내놓고, 인접한 node들 사이에서 boundary에
공유된 색의 slot들을 동일시한다. 그 동일시를 한 걸음씩이 아니라
``\lean{colorGraph}에서의 도달 가능성''으로 한꺼번에 정의했기 때문에
동치관계임 (\lean{occRel\_refl/symm/trans})이 곧바로 나온다. 한 걸음씩
정의하고 동치폐포를 취하는 방식보다 다루기 쉽다.

Transitivity 증명에서 ``중간 slot의 색이 같아야 한다''는 관찰
(\lean{hbi.symm.trans hbj})이 필요한데, 이는 slot이 색을 결정하기 때문이다.
\end{remark}

주요 결과:
\begin{itemize}
  \item \lean{decodedColor\_isBagColoring}: 복원된 coloring이 bag-injective하다.
  \item \lean{decode\_encode\_letter\_equivalent}: decode한 뒤 다시 encode하면
    원래 letter와 \emph{동치}다 (\lean{SigmaLetter.Equivalent}, 색을 통한
    동일시).
  \item \lean{decode\_encode\_iso}: encode한 뒤 decode하면 원래 $\tau_P$-graph와
    \emph{isomorphic}하다 (\lean{τPGraph.Iso}).
\end{itemize}

\begin{remark}[두 방향의 ``역''이 서로 다른 강도다]
\label{mso:rem:twoinverses}
Encode-after-decode는 letter의 \emph{동치}로, decode-after-encode는 graph의
\emph{isomorphism}으로 서술된다. 등호가 아닌 이유는 quotient로 만든 vertex가
원래 vertex와 같은 타입이 아니기 때문이다. 이 비대칭은 형식화의 한계가 아니라
수학적으로 올바른 진술이다.

Decode-after-encode 증명의 뼈대는 \lean{occurrenceOriginal}: encoding의 각
slot $(t,i)$에 대해 ``$t$의 bag에서 색 $i$를 갖는 원래 vertex''를 돌려주는
함수를 만들고, 이것이 \lean{OccRel}에 대해 well-defined이며
(\lean{occRel\_iff\_occurrenceOriginal\_eq}) bijective임
(\lean{decodeToOriginal\_bijective}) 을 보인다. 여기서 다시 bag-injectivity가
핵심이다.
\end{remark}

이 절 전체는 Courcelle 정리의 \emph{증명 경로에는 필요하지 않다}. 필요한 것은
\S\ref{mso:thm:encodeislegal}의 정방향뿐이다. Decoding은 ``legal $\Sigma$-tree와
colored decomposition이 본질적으로 같은 것''이라는 참조 문서의 주장을 완전히
형식화하기 위해 있으며, automata 쪽에서 언어의 비어 있지 않음을 다룰 때
의미를 준다.

\section{Tree model로 가는 다리}
\label{mso:sec:treemodel}

\emph{파일: \leanfile{GraphMSO/Decomp/treeModel.lean}}

\begin{definition}[\lean{SigmaTree.toTreeModel}]
\label{mso:def:totreemodel}
$\Sigma$-tree를 \leanfile{Lean\_Basics.tex}의 tree language model
(\lean{basics:def:treemodel})로 본다: node는 그대로, \lean{parentRel}은
\lean{IsChild}, \lean{label}은 \lean{letter}.
\end{definition}

짧지만 결정적인 파일이다. 두 개의 등식이 이후 모든 계산의 기반이 된다.

\begin{theorem}
\label{mso:thm:modelbridge}
Encoding에 대해:
\begin{itemize}
  \item \lean{encode\_toTreeModel\_parentRel}: model의 parent 관계가
    decomposition의 \lean{IsChild}와 \emph{정의상} 같다 (증명이 \lean{rfl}).
  \item \lean{encode\_toTreeModel\_graph}: model의 대칭화된 parent graph
    (\leanfile{Lean\_Basics.tex}의 \lean{basics:def:treemodelgraph})가
    \emph{decomposition tree 자체}와 같다.
\end{itemize}
\end{theorem}

\begin{remark}[왜 이 두 등식이 중요한가]
\leanfile{Lean\_Basics.tex}에서 tree language의 유도 formula들
(\lean{conn}, \lean{top}, \lean{dangle})의 의미론은 ``\lean{TreeModel.graph}에서의
연결성''이나 ``\lean{parentRel}''로 서술되어 있었다. 위 두 등식이 있으면
그 진술들이 encoding 위에서 곧바로 \emph{decomposition tree에서의 연결성}과
\emph{decomposition의 child 관계}에 대한 진술로 바뀐다.

두 번째 등식은 \leanfile{Lean\_Basics.tex}의
\lean{adj\_iff\_isChild\_or\_isChild}로 증명된다. 즉 ``tree adjacency는
두 방향 중 하나의 child 관계''라는 사실이 여기서 값을 한다.
\end{remark}

\section{(3a) Defining pair와 defining tuple}
\label{mso:sec:definingpairs}

\emph{파일: \leanfile{GraphMSO/Decomp/definingPairs.lean}}

\emph{참조: \leanfile{lecture\_note\_expanded.tex} \S``Combinatorial properties
of defining tuples'': Lemma \lean{lem:defining-pairs},
\lean{lem:distinct-pairs}, \lean{lem:edges-bags}, \lean{lem:defining-tuples}}

이 절은 순수 조합론이다. Formula는 전혀 나오지 않는다. 그러나 \emph{formula가
말하게 될 내용}을 미리 조합론적으로 확립해 둔다.

\subsection{Defining pair}

핵심 착상은 이렇다. Tree MSO는 원래 graph의 vertex를 직접 언급할 수 없다.
언급할 수 있는 것은 tree의 node 집합뿐이다. 그런데 vertex $v$는 쌍
\[
  (\BAGS(v),\ c(v)) \quad (\text{node 집합, 색})
\]
로 완전히 결정된다. 이 쌍을 $v$의 \emph{defining pair}라 부르고, ``어떤 쌍이
defining pair인가''를 tree MSO로 표현할 수 있게 만드는 것이 목표다.

\begin{definition}[\lean{IsDefiningPair}]
\label{mso:def:isdefiningpair}
node 집합 $U$와 색 $i$에 대한 다섯(코드에서는 여섯 항목) 조건:
\begin{enumerate}
  \item $U$가 비어 있지 않다;
  \item $U$가 preconnected하다;
  \item $U$의 모든 node의 bag에 색 $i$인 vertex가 있다;
  \item $U$의 원소 중 \emph{부모가 $U$ 안에 있는} 것은 그 adhesion에 색 $i$인
    vertex가 있다;
  \item $U$의 원소 중 \emph{부모가 $U$ 밖에 있는} 것은 그 adhesion에 색 $i$인
    vertex가 없다;
  \item $U$ 밖이면서 \emph{부모가 $U$ 안에 있는} node는 그 adhesion에 색 $i$인
    vertex가 없다.
\end{enumerate}
\end{definition}

\begin{remark}[``top node''를 말하지 않는다]
\label{mso:rem:topbyparent}
조건 4--6에 주목하라. 참조 문서는 ``$U$의 top node''와 ``$U$에 매달린
자식''을 명시적으로 언급하지만, Lean 정의는 그 대신 \emph{부모가 집합을
벗어나는가}로 표현한다.

이것이 결정적인 설계 선택이다. Tree MSO의 formula \lean{top}과 \lean{dangle}
(\leanfile{Lean\_Basics.tex}의 \lean{basics:def:treederived})이 정확히 그
모양이기 때문이다. Preconnected한 $U$에 대해 ``부모가 $U$ 밖''은 ``top
node''와 동치이지만 (\leanfile{Lean\_Basics.tex}의
\lean{basics:thm:topunique}), formula 쪽 모양에 미리 맞춰 두면
\S\ref{mso:sec:recognition}의 대응 증명이 거의 기계적이 된다. 만약 top node로
정의했다면 매번 그 동치를 오가야 했을 것이다.

조건 6의 ``$U$ 밖이지만 부모가 $U$ 안''이 \lean{dangle}이다.
\end{remark}

\begin{theorem}[\lean{isDefiningPair\_iff}, 참조 문서 Lemma \lean{lem:defining-pairs}]
\label{mso:thm:definingpairiff}
$c$가 bag-injective이면
\[
  \lean{IsDefiningPair}\ c\ U\ i \iff \exists v.\ c(v) = i \land \BAGS(v) = U.
\]
\end{theorem}

\begin{remark}[역방향 증명 방식]
\label{mso:rem:definingpairproof}
정방향 (\lean{isDefiningPair\_BAGS}) 은 각 조건을 확인하면 되므로 쉽다.
역방향 (\lean{exists\_of\_isDefiningPair}) 이 본론이며 세 단계다.

\emph{1단계.} $U$의 top node에서 조건 3으로 색 $i$인 vertex $v$를 얻는다.
그리고 \emph{$U$의 모든 node가 $v$를 담는다}는 것을 top node에서 내려가는
ancestor chain에 대한 귀납으로 보인다. 귀납 단계에서 조건 4가 자식의 adhesion에
색 $i$인 vertex를 주고, bag-injectivity가 그것이 $v$와 같음을 강제한다.
\emph{즉 조건 4가 ``$v$가 계속 살아 있다''를 보장한다.}

\emph{2단계.} $U$의 top node가 $\BAGS(v)$의 top node와 같음을 보인다. 만약
$\BAGS(v)$가 더 위로 올라간다면 $U$의 top node의 adhesion에 $v$가 들어가는데,
이는 조건 5에 위배된다. \emph{즉 조건 5가 ``$v$가 위로 새지 않는다''를
보장한다.}

\emph{3단계.} 역으로 $\BAGS(v)$의 모든 node가 $U$에 속함을 같은 방식의 귀납으로
보인다. 어떤 node가 $U$를 벗어나면서 $v$를 담으면 그 node는 dangling이고
adhesion에 $v$가 있으므로 조건 6에 위배된다. \emph{즉 조건 6이 ``$v$가 아래로
새지 않는다''를 보장한다.}

세 조건 4, 5, 6이 각각 ``유지'', ``위쪽 차단'', ``아래쪽 차단''을 담당한다는
것이 이 정의의 요점이다. 두 귀납 모두 \leanfile{Lean\_Basics.tex}의 볼록성
정리 (\lean{mem\_of\_isAncestor\_of\_isAncestor},
\lean{mem\_BAGS\_of\_isAncestor\_of\_isAncestor}) 를 사용한다.
\end{remark}

\subsection{Distinctness와 atom 읽기}

참조 문서 Lemma \lean{lem:distinct-pairs}에 대응:
\lean{disjoint\_BAGS\_of\_ne\_of\_color\_eq} (같은 색의 서로 다른 vertex는
$\BAGS$가 서로소), \lean{not\_adj\_of\_ne\_of\_color\_eq} (인접하지 않음),
\lean{eq\_of\_BAGS\_eq\_of\_color\_eq} (vertex는 defining pair로 결정됨).
셋 다 bag-injectivity의 직접적 귀결이다.

참조 문서 Lemma \lean{lem:edges-bags}에 대응하는 것이 다음 둘이며, 이들이
atomic formula의 번역을 정당화한다.

\begin{theorem}[\lean{adj\_iff\_exists\_mem\_BAGS\_adjOnColors}]
\label{mso:thm:adjfrombags}
\[
  u \sim_G v \iff \exists t.\ t \in \BAGS(u) \land t \in \BAGS(v) \land
    \lean{encodeLetter}(t).\lean{AdjOnColors}\ c(u)\ c(v).
\]
\end{theorem}

\begin{theorem}[\lean{vpred\_iff\_exists\_mem\_BAGS\_tagOnColor}]
\label{mso:thm:predfrombags}
\[
  \lean{vpred}\ p\ v \iff \exists t \in \BAGS(v).\
    \lean{encodeLetter}(t).\lean{TagOnColor}\ p\ c(v).
\]
\end{theorem}

즉 원래 graph의 adjacency와 unary predicate를 \emph{letter만 보고} 판정할 수
있다. 정방향은 edge coverage 공리로 공통 bag을 찾고, 역방향은
bag-injectivity로 letter의 존재 한정을 실제 vertex로 되돌린다.
(\lean{vpred} 쪽은 사실 ``어떤 $t$''를 ``모든 $t$''로 강화할 수 있다:
\lean{encodeLetter\_tagOnColor\_iff\_vpred}.)

\subsection{Defining tuple}

Vertex 하나가 아니라 vertex \emph{집합}을 표현하려면 좌표가 여러 개 필요하다.

\begin{definition}[\lean{definingTuple}]
\label{mso:def:definingtuple}
vertex 집합 $S$에 대해
\[
  \lean{definingTuple}\ c\ S : \lean{Fin}\ k \to \lean{Set}\ \lean{Node},
  \qquad
  i \mapsto \bigcup_{v \in S,\ c(v) = i} \BAGS(v).
\]
즉 좌표 $i$는 $S$의 색-$i$ 원소들의 $\BAGS$를 모두 모은 것이다.
\end{definition}

\begin{remark}[왜 좌표가 $\omega+1$개 필요한가]
같은 색의 서로 다른 vertex들은 $\BAGS$가 서로소이므로
(\lean{disjoint\_BAGS\_of\_ne\_of\_color\_eq}), 한 좌표 안에서 합집합을 취해도
정보가 섞이지 않는다 --- 각 연결 성분이 하나의 vertex에 해당한다. 반면 색이
다른 vertex들은 $\BAGS$가 겹칠 수 있으므로 좌표를 나눠야 한다. 그래서 색의
개수 $\omega+1$개의 좌표가 필요하고, 이것이 \S\ref{mso:sec:translation}에서
graph 변수 하나가 tree 변수 $\omega+1$개로 확장되는 이유다.
\end{remark}

관련 정리들:
\begin{itemize}
  \item \lean{definingTuple\_singleton}: 한 점 집합의 tuple은 그 vertex의 색
    좌표에서 $\BAGS(v)$, 나머지 좌표에서 $\emptyset$.
  \item \lean{eq\_of\_definingTuple\_singleton\_eq}: vertex는 자기 tuple로
    결정된다 (등호 atom의 번역 근거).
  \item \lean{mem\_iff\_BAGS\_subset\_definingTuple}:
    $v \in S \iff \BAGS(v) \subseteq \lean{definingTuple}\ S\ (c(v))$
    (소속 atom의 번역 근거).
  \item \lean{exists\_definingTuple\_eq\_iff}: 주어진 tuple이 어떤 vertex
    집합에서 오는지의 판정 --- ``각 좌표의 모든 node가, 그 좌표에 포함되는
    defining pair로 덮인다''. 이것이 set-recognition formula의 모양 그대로다.
    (참조 문서 Lemma \lean{lem:defining-tuples}.)
\end{itemize}

\section{(3b) Recognition formula}
\label{mso:sec:recognition}

\emph{파일: \leanfile{GraphMSO/Decomp/recognition.lean}}

\emph{참조: \leanfile{lecture\_note\_expanded.tex} Lemma \lean{lem:vtx-i},
\lean{lem:vtx}, \lean{lem:set}, \lean{lem:adj}, \lean{lem:pred}, \lean{lem:eq},
\lean{lem:cont}, Corollary \lean{cor:recognition-enc}}

이제 \S\ref{mso:sec:definingpairs}의 조합론적 사실을
\leanfile{Lean\_Basics.tex}에서 \emph{syntax로만} 정의해 둔 formula들
(\lean{basics:def:recognitionformulas})과 연결한다. 이 절의 정리들은 모두
같은 모양이다: ``formula가 성립한다 $\iff$ 조합론적 조건이 성립한다''.

\begin{definition}[\lean{legalFormula}]
\label{mso:def:legalformula}
\[
  \lean{legalFormula}\ P\ \omega :=
    \lean{Formula.legal}\ \{a \mid a.\lean{RootEmpty}\}\
      \{(q_1,q_2) \mid \lean{Compatible}\ q_1\ q_2\}.
\]
\end{definition}

\begin{theorem}[\lean{satisfiesAt\_legalFormula\_iff}]
\label{mso:thm:legalformula}
임의의 $\Sigma$-tree $S$에 대해 $S \models \lean{legalFormula} \iff S.\lean{IsLegal}$.
\end{theorem}

\begin{remark}
참조 문서는 이 formula를 letter마다의 $P_a$에 대한 이중 유한 논리합으로
쓰지만, Lean에서는 letter \emph{집합}을 원자식에 직접 넣는
\lean{labelMem}/\lean{labelMem₂} 덕분에 (\leanfile{Lean\_Basics.tex}의
\lean{basics:rem:treesimplify}) formula가 상수 크기다. 증명도 그만큼 짧아서,
``parent가 없는 node는 정확히 root''(\lean{forall\_not\_isChild\_iff})만
확인하면 된다.
\end{remark}

\begin{theorem}[\lean{satisfiesAt\_definingPair\_iff}, 참조 문서 Lemma \lean{lem:vtx-i}]
\label{mso:thm:definingpairformula}
Encoding 위에서, formula
$\lean{definingPair}\ \{a \mid a.\lean{HasVertex}\ i\}\ \{a \mid a.\lean{RootContains}\ i\}\ Z$
가 성립하는 것은 $(\rho(Z), i)$가 defining pair인 것과 동치다.
\end{theorem}

\begin{remark}[증명이 왜 기계적인가]
\label{mso:rem:recognitionmechanical}
증명은 네 개의 private 보조정리 (\lean{comp\_A\_iff}, \lean{comp\_Rnt\_iff},
\lean{comp\_Rtop\_iff}, \lean{comp\_Rdang\_iff}) 로 나뉘며, 각각이
\lean{IsDefiningPair}의 조건 3, 4, 5, 6에 대응한다. 각 보조정리는
\leanfile{Lean\_Basics.tex}의 semantic characterization
(\lean{satisfiesAt\_top\_iff} 등)과 \S\ref{mso:thm:modelbridge}의 model bridge
등식, 그리고 \S\ref{mso:def:encodeletter}의 letter 읽기 보조정리
(\lean{encodeLetter\_rootContains\_iff} 등)를 차례로 \lean{rw}로 적용한
것이다.

세 층 (formula semantics / model bridge / encoding) 을 미리 각각 완결해
두었기 때문에 이 지점에서는 재작성만 하면 된다. 마지막에 남는 것은
\lean{top} formula의 ``root이거나 부모가 밖에 있다'' 형태와
\lean{IsDefiningPair}의 ``부모가 밖에 있다'' 형태 사이의 사소한 case 분석뿐인데,
root의 adhesion이 $\emptyset$이라는 사실이 그 간극을 메운다.

\S\ref{mso:rem:topbyparent}에서 예고한 설계의 배당금이 여기서 나온다.
\end{remark}

이어지는 정리들:
\begin{itemize}
  \item \lean{satisfiesAt\_vtxTuple\_iff} (Lemma \lean{lem:vtx}): tuple이 어떤
    vertex의 defining tuple이다. 증명은 ``어떤 좌표가 defining pair이고 나머지
    좌표는 비어 있다''를 \S\ref{mso:thm:definingpairiff}로 vertex로 바꾼다.
  \item \lean{satisfiesAt\_setTuple\_iff} (Lemma \lean{lem:set}): tuple이 어떤
    vertex 집합의 defining tuple이다. \lean{exists\_definingTuple\_eq\_iff}가
    그대로 대응한다.
  \item \lean{satisfiesAt\_adjTuple\_iff} (Lemma \lean{lem:adj}),
    \lean{satisfiesAt\_predTuple\_iff} (Lemma \lean{lem:pred}),
    \lean{satisfiesAt\_eqTuple\_iff} (Lemma \lean{lem:eq}),
    \lean{satisfiesAt\_contTuple\_iff} (Lemma \lean{lem:cont}): 네 atomic
    formula가 각각 \S\ref{mso:thm:adjfrombags}, \S\ref{mso:thm:predfrombags},
    \lean{eq\_of\_definingTuple\_singleton\_eq},
    \lean{mem\_iff\_BAGS\_subset\_definingTuple}에 대응한다.
\end{itemize}

이 목록 전체가 참조 문서의 Corollary \lean{cor:recognition-enc}다.

\section{(3c) Translation \texorpdfstring{$\theta^\ast$}{theta star}}
\label{mso:sec:translation}

\emph{파일: \leanfile{GraphMSO/Decomp/translation.lean}}

\emph{참조: \leanfile{lecture\_note\_expanded.tex} \S``The formula
translation'', Lemma \lean{lem:translation}}

\subsection{변수 할당}

Graph 변수 하나가 tree 변수 $\omega+1$개로 확장되므로, 충돌 없는 할당 규칙이
필요하다.

\begin{definition}[\lean{fvBlock}, \lean{svBlock}]
\label{mso:def:blocks}
\[
  \lean{fvBlock}\ \omega\ x\ i := (\omega+1)(2x) + i,
  \qquad
  \lean{svBlock}\ \omega\ X\ i := (\omega+1)(2X+1) + i
\]
($i \in \lean{Fin}(\omega+1)$). \lean{blockList}는 한 block을 list로 만든 것으로,
block quantifier에 넘긴다.
\end{definition}

\begin{remark}[$(\omega+1)$진법 자리표기]
\label{mso:rem:positional}
이 할당은 $(\omega+1)$진법 자리표기다: 몫이 graph 변수 번호를 (parity와 함께)
주고, 나머지가 block 안의 좌표를 준다. 핵심 보조정리
\lean{block\_eq\_iff}가 그 사실을 나눗셈으로 증명한다:
\[
  k a + i = k b + j \iff a = b \land i = j \qquad (i, j < k).
\]
여기서 $\lean{Fin}\ k$의 원소가 $k$ 미만이라는 것 (\lean{Fin.isLt}) 이 쓰인다.

그 위에 필요한 모든 injectivity와 disjointness가 따라 나온다:
\lean{fvBlock\_injective}, \lean{svBlock\_injective},
\lean{fvBlock\_ne\_svBlock} (parity가 다르므로 --- $2x$와 $2X+1$),
\lean{fvBlock\_ne\_fvBlock}, \lean{svBlock\_ne\_svBlock}, 그리고 네 개의
``$\notin \lean{blockList}$'' 보조정리.

이 보조정리들이 왜 이렇게 많이 필요한가? \leanfile{Lean\_Basics.tex}의
\lean{basics:rem:setblock}에서 설명한 대로 \lean{Assignment.setBlock}의 핵심
등식이 block 열거의 injectivity를 요구하고, quantifier 경우마다 ``이 block을
바꿔도 다른 block은 그대로''를 보여야 하기 때문이다.
\end{remark}

\subsection{Translation}

\begin{definition}[\lean{Formula.translate}]
\label{mso:def:translate}
$\tau_P$ formula를 $\Sigma_\omega$ 위의 tree formula로 옮긴다.
\begin{itemize}
  \item Atom은 \S\ref{mso:sec:recognition}의 tuple formula로:
    \lean{equal} $\mapsto$ \lean{eqTuple}, \lean{adj} $\mapsto$ \lean{adjTuple},
    \lean{pred} $\mapsto$ \lean{predTuple}, \lean{inSet} $\mapsto$
    \lean{contTuple}. Letter 집합 인자는 각각
    $\{a \mid a.\lean{HasVertex}\ i\}$, $\{a \mid a.\lean{RootContains}\ i\}$,
    $\{a \mid a.\lean{AdjOnColors}\ i\ j\}$, $\{a \mid a.\lean{TagOnColor}\ p\ i\}$.
  \item 논리 결합자는 그대로 통과한다.
  \item Quantifier는 \emph{guard된 block quantifier}가 된다:
    \[
      \exists x\, \varphi \;\mapsto\;
        \exists \lean{fvBlock}(x).\ \bigl(\lean{vtxTuple}(\lean{fvBlock}\ x) \land \varphi^\ast\bigr),
    \]
    \[
      \forall x\, \varphi \;\mapsto\;
        \forall \lean{fvBlock}(x).\ \bigl(\lean{vtxTuple}(\lean{fvBlock}\ x) \to \varphi^\ast\bigr),
    \]
    집합 quantifier는 \lean{svBlock}과 \lean{setTuple}로 같은 모양.
\end{itemize}
\end{definition}

\begin{remark}[전칭 quantifier를 직접 번역한다]
$\forall$을 $\lnot\exists\lnot$로 우회하지 않고 직접 guarded universal block으로
번역한 점을 짚어 둔다. 이것이 가능한 이유는 tree language가 $\forall$과
$\to$를 primitive로 갖기 때문이다 (\leanfile{Lean\_Basics.tex}의 설계).
결과적으로 번역된 formula의 크기가 원래 formula의 구조를 그대로 반영하고,
\leanfile{Lean\_Automata.tex}의 compile 단계에서 $\forall$ 경우를 따로
처리할 수 있게 된다.

Guard가 필요한 이유는 명백하다: tree 변수 block에 아무 node 집합이나 넣을 수
있지만, 그중 실제로 graph 변수에 대응하는 것은 defining tuple인 것들뿐이다.
Guard가 그 제약을 건다.
\end{remark}

\subsection{정확성}

\begin{theorem}[\lean{satisfiesAt\_translate\_iff}, 참조 문서 Lemma \lean{lem:translation}]
\label{mso:thm:translatecorrect}
Encoding 위에서, $\tau_P$ formula $\theta$와 graph assignment $\alpha$,
tree assignment $\rho$에 대해 다음을 가정하자.
\begin{itemize}
  \item $\theta$의 모든 자유 first-order 변수 $x$에 대해, 어떤 $v$가 있어
    $\alpha(x) = \lean{some}\ v$이고 모든 $i$에 대해
    $\rho(\lean{fvBlock}\ x\ i) = \lean{definingTuple}\ \{v\}\ i$;
  \item $\theta$의 모든 자유 second-order 변수 $Y$에 대해, 모든 $i$에 대해
    $\rho(\lean{svBlock}\ Y\ i) = \lean{definingTuple}\ (\alpha(Y))\ i$.
\end{itemize}
그러면
\[
  \lean{encode}.\lean{toTreeModel},\ \rho \models \theta^\ast
  \iff
  (V, G, \lean{vpred}),\ \alpha \models \theta.
\]
\end{theorem}

\begin{remark}[가정의 형태가 증명 설계의 전부다]
\label{mso:rem:translationhyp}
이 정리에서 배울 점은 결론이 아니라 \emph{가정의 모양}이다.

순진하게 진술하면 ``닫힌 formula에 대해 $G \models \theta \iff
\lean{encode} \models \theta^\ast$''로 충분해 보인다. 그러나 이것을 formula에
대한 귀납으로 증명하려면 quantifier 경우에서 자유 변수가 생기므로, 열린
formula에 대한 진술이 필요하다. 그리고 열린 formula에 대해 서술하려면
``graph assignment와 tree assignment가 대응한다''를 말해야 한다.

여기서 선택이 갈린다. 대응을 ``모든 변수에서''로 요구하면 quantifier 경우에서
새로 묶인 변수에 대해서도 대응을 만들어야 하는데, 그 변수는 원래 assignment에
없던 것이라 곤란해진다. 그래서 이 형식화는 대응을 \emph{$\theta$의 자유
변수에서만} 요구한다. 그러면 $\exists x\,\varphi$의 경우:
\begin{itemize}
  \item ($\Leftarrow$) graph 쪽 witness $v$가 주어지면 tree 쪽 witness로
    $\lean{definingTuple}\ \{v\}$를 \lean{setBlock}으로 넣는다. Guard는
    \S\ref{mso:sec:recognition}의 \lean{satisfiesAt\_vtxTuple\_iff}가 만족시켜
    주고, 나머지 block이 안 건드려졌다는 것은
    \lean{setBlock\_so\_of\_forall\_ne}와 \S\ref{mso:rem:positional}의
    disjointness 보조정리가 준다.
  \item ($\Rightarrow$) tree 쪽 witness가 guard를 만족하므로
    \lean{satisfiesAt\_vtxTuple\_iff}가 그것이 어떤 $v$의 defining tuple임을
    준다. 그 $v$가 graph 쪽 witness다.
\end{itemize}
어느 방향이든 귀납 가설의 가정이 정확히 재생산된다. 파일 주석의 표현대로
``self-sustaining''하며, 그 덕분에 tree language에 대한 별도의 coincidence
lemma나 capture-avoidance 장치를 만들 필요가 없다.

이것이 \leanfile{Lean\_Basics.tex}에서 tree language의 자유 변수 개념을
아예 정의하지 않은 이유이기도 하다. 필요가 없었던 것이다.
\end{remark}

\begin{corollary}[\lean{satisfies\_legalFormula\_conj\_translate\_iff}]
\label{mso:cor:finaldisplay}
$\theta$가 sentence ($\lean{freeFO} = \lean{freeSO} = \emptyset$) 이면
\[
  \lean{encode}.\lean{toTreeModel} \models
    \bigl(\varphi_{\mathrm{legal}} \land \theta^\ast\bigr)
  \iff
  (V, G, \lean{vpred}) \models \theta.
\]
\end{corollary}

이것이 참조 문서의 최종 display이며, 이 문서의 목표다. 왼쪽은 순수하게
$\Sigma$-tree에 대한 tree MSO sentence이므로, 이제
\leanfile{Lean\_Automata.tex}가 그것을 tree automaton으로 컴파일할 수 있다.

\begin{remark}[$\varphi_{\mathrm{legal}}$이 왜 붙는가]
Translation $\theta^\ast$ 혼자로는 부족하다. 임의의 $\Sigma$-tree --- legal하지
않은 것도 포함 --- 위에서 $\theta^\ast$가 무엇을 뜻하는지는 통제되지 않기
때문이다. $\varphi_{\mathrm{legal}}$을 붙여 ``실제로 어떤 decomposition의
encoding인 tree''로 범위를 제한해야 한다.

\S\ref{mso:def:recognitionformulas}에서 언급했듯, \lean{definingPair} formula
자체에는 legality 연언항을 넣지 않았다. Encoding 위에서는 자동으로 성립하고,
최종 sentence에서 최상위에 한 번만 붙이는 편이 formula 크기 면에서 유리하기
때문이다.
\end{remark}

\section{요약: 참조 문서와의 대응}
\label{mso:sec:summary}

\begin{center}
\begin{tabular}{p{0.44\textwidth}p{0.48\textwidth}}
\hline
\leanfile{lecture\_note\_expanded.tex} & Lean \\
\hline
Corollary ($\MSOtwo$ properties) &
  \lean{incidence\_reduction} (\S\ref{mso:cor:incidencereduction}) \\
\S Rooted graphs and gluing &
  \lean{KRootedGraph}, \lean{GluingData}, \lean{MultiGluingData} \\
Lemma \lean{lem:cone-struct} &
  \lean{nodeConeGraph.lean}의 구조 보조정리들 \\
Lemma \lean{lem:cone-gluing} &
  \lean{cone\_isMultiGluing} \\
Definition \lean{def:sigma-tree} &
  \lean{SigmaLetter}, \lean{SigmaTree} (단, unordered --- \S\ref{mso:rem:unordered}) \\
Definition (legal) &
  \lean{SigmaLetter.Compatible}, \lean{SigmaTree.IsLegal} \\
$\varphi_{\mathrm{legal}}$ lemma &
  \lean{satisfiesAt\_legalFormula\_iff} \\
\S Decoding a legal tree, Lemma \lean{lem:decode-enc} &
  \lean{decoding.lean} (\lean{decode\_encode\_iso} 등) \\
Lemma \lean{lem:legal-iff-enc} (정방향) &
  \lean{encode\_isLegal} \\
Lemma \lean{lem:defining-pairs} &
  \lean{isDefiningPair\_iff} \\
Lemma \lean{lem:distinct-pairs} &
  \lean{disjoint\_BAGS\_of\_ne\_of\_color\_eq} 등 \\
Lemma \lean{lem:edges-bags} &
  \lean{adj\_iff\_exists\_mem\_BAGS\_adjOnColors} \\
Lemma \lean{lem:defining-tuples} &
  \lean{exists\_definingTuple\_eq\_iff} \\
Lemma \lean{lem:vtx-i} / \lean{lem:vtx} / \lean{lem:set} &
  \lean{satisfiesAt\_definingPair\_iff} / \lean{\_vtxTuple\_iff} / \lean{\_setTuple\_iff} \\
Lemma \lean{lem:adj}/\lean{lem:pred}/\lean{lem:eq}/\lean{lem:cont} &
  \lean{satisfiesAt\_adjTuple\_iff} 외 3 \\
Corollary \lean{cor:recognition-enc} &
  위 일곱 정리의 묶음 \\
\S The formula translation &
  \lean{Formula.translate} \\
Lemma \lean{lem:translation} &
  \lean{satisfiesAt\_translate\_iff} \\
최종 display &
  \lean{satisfies\_legalFormula\_conj\_translate\_iff} \\
\hline
\end{tabular}
\end{center}

\begin{remark}[다음 단계]
남은 것은 참조 문서의 Theorem \lean{thm:automata} (effective
$\MSO$--tree-automaton correspondence)이며, 그것이
\leanfile{Lean\_Automata.tex}의 내용이다. 거기서 \S\ref{mso:rem:unordered}의
unordered/ordered 간극이 메워지고, \S\ref{mso:rem:lettercolors}의 alphabet
유한성이 비로소 쓰인다.
\end{remark}

\end{document}
